% Complete analytic module; uses HKO coordinates and thm:mc-derivative.
% Four literal family predicates are checked by certify_amplification.py.
\subsection{A first-order obstruction for two characteristic-invariant cuts}
\label{subsec:ha-two-hinge}

The seven-cut threshold above concerns ten fixed radial profiles. A different
obstruction applies to \emph{any two} affine hinges that are constant along
their own facet characteristic directions. Affine offsets, dependence on
the facet-edge coordinate, and unequal cut speeds are all allowed.

For a nonnegative function $f$ on the labeled facet union, write
$R(f)=\sup_\gamma\int f\,d\mu_\gamma$, with the supremum over all finite
maximizing words, or equivalently their weak occupation limits.

\begin{theorem}[Two invariant hinges]\label{thm:ha-two-hinge}
Suppose $f_i(x)=(D_i\cdot x)_+$ on a single HKO facet $F_i$, vanishing on
the other labeled facets, and $D_i\cdot Ja_i=0$. Then
\begin{equation}\label{eq:ha-amplification}
                         R(f_i)\geq2\nu(f_i).
\end{equation}
Consequently, adding any two such cuts while retaining every old inequality
has nonpositive first systolic variation. The two cuts may concern distinct
facets or the same facet.
\end{theorem}

This is a first-order statement for characteristic-invariant affine hinges.
It does not prove local maximality among twelve-facet bodies, decide a
zero-derivative direction at higher order, or cover non-invariant hinges.

\paragraph{Quotient coordinates and volume.}
The facet-transitive symmetries used in Proposition~\ref{prop:ha-weighted}
reduce the proof to the q-facet labeled zero. Put
\[
 r=\sqrt5,\quad N=r-1,\quad e=J_2n_0/N,\quad
 m=\frac{1-\cos(\pi/5)}2,\quad h=\frac{1+\cos(\pi/5)}2,
\]
and on that facet use
\[
 \xi=\frac{e\cdot q}{\sin(\pi/5)},\qquad
 \eta=\frac{e\cdot p-m}{h}.
\]
The quotient ranges over $[-1,1]^2$. Every affine function invariant along
$Ja_0$ is $\alpha\xi+\beta\eta+\gamma$: its normal coordinate is constant
on the facet, and invariance excludes precisely the remaining partner
coordinate. Under facet-uniform measure, $\xi$ is uniform on $[-1,1]$ and
independent of $\eta$. Put
\[
 k=r-2,\qquad a=\frac{5-r}{20},\quad b=\frac r{10},\qquad
 A=\frac a2,\quad B=\frac{5+r}{40},\quad C=\frac{5+2r}{10}.
\]
Projection of the partner pentagon gives the probability density
\begin{equation}\label{eq:ha-quotient-density}
 p(y)=\begin{cases}b(y+3),&-1\leq y\leq k,\\
                   C(1-y),&k\leq y\leq1,
       \end{cases}
 \qquad \mu:=\int yp(y)\,dy=-1+\frac{2r}{5}.
\end{equation}
Indeed its vertex projection levels are $-1,k,1$, and the transverse slice
width is affine between them. The exact projection, normalization and mean
are also checked in the accompanying four-family certificate. Each facet
has cone-volume mass $1/10$, so
\begin{equation}\label{eq:ha-quotient-volume}
 \nu(f_0)=\frac1{10}\mathbb E
                (\alpha\xi+\beta\eta+\gamma)_+.
\end{equation}

\paragraph{Four literal shortest-word families.}
The following families realize each edge coordinate $\xi=\pm1$ and every
listed value of $\eta=y$. The facet-zero occupation weight is
\begin{equation}\label{eq:ha-two-weight-curves}
 \beta_L(y)=A(y+3)\quad(-1\leq y\leq k),\qquad
 \beta_U(y)=B(1-y)\quad(k\leq y\leq1).
\end{equation}
At the common endpoint, $\beta_L(k)=b$ whereas $\beta_U(k)=a$; the two
families need not join with the same weight.
Here are their complete ordered words and parameter intervals:
\[
\begin{array}{c|c|c|l}
 \text{family}&\xi& y(z),\quad0\leq z\leq1&\sigma\\ \hline
 L_+& 1&-1+(k+1)z&(0,1,7,6,3,2,9)\\
 L_-&-1&-1+(k+1)z&(0,6,3,2,9,8,4)\\
 U_+& 1& k+(1-k)z&(0,7,3,9,4,5,1)\\
 U_-&-1& k+(1-k)z&(0,4,6,5,2,1,8)
\end{array}
\]
In each listed order, the full weight vector is
$\beta(z)=(1-z)\beta^-+z\beta^+$, with
\[
\begin{array}{c|c|c}
 &\beta^-&\beta^+\\ \hline
 L_+&(a,a,a,a,b,0,b)&(b,0,a,a,a,a,b)\\
 L_-&(a,b,0,b,a,a,a)&(b,b,a,a,a,a,0)\\
 U_+&(a,b,b,a,0,a,a)&(0,b,a,a,a,a,b)\\
 U_-&(a,a,a,a,b,0,b)&(0,b,a,a,a,a,b)
\end{array}
\]
These are catalogue charts 115, 148, 72 and 19, with their parameters oriented
to make $y$ increase. Direct exact arithmetic checks closure, nonnegative
weights and $Q=\tau/5$ at both endpoints and the midpoint. The KKT
reconstruction checks all facet inequalities at all loop vertices,
$\xi=\pm1$, the displayed affine $y$ coordinate, and
\eqref{eq:ha-two-weight-curves}. Endpoint checks extend to the whole family:
the weights and reconstructed positions are affine, and the objective is
quadratic. Interior weights are positive. Only the existence of these four
shortest families is needed below; no numerical search result is a premise.

\begin{lemma}[Scalar hinge bound]\label{lem:ha-scalar-amplification}
For $g(y)=(uy+v)_+$, set
\[
 I(g)=\int_{-1}^1g(y)p(y)\,dy,\qquad
 M(g)=\max\left\{\sup_{[-1,k]}A(y+3)g(y),\
                       \sup_{[k,1]}B(1-y)g(y)\right\}.
\]
Then $I(g)\leq5M(g)$.
\end{lemma}

\begin{proof}
Constants follow from $M(g)\geq bg$ and $5b=r/2>1$.
Positive homogeneity reduces the other cases to decreasing hinges
$(z-y)_+$ and increasing hinges $(y-z)_+$, for real thresholds $z$.

For a decreasing hinge with $z\leq-1$ the claim is trivial. On $[-1,1]$,
convexity in $z$ bounds its integral by the endpoint secant:
\[
 I(z)\leq\frac{z+1}{2}(1-\mu),\qquad M(z)\geq2A(z+1).
\]
The required comparison follows from
$20A-(1-\mu)=(5-r)/10>0$.
For $z\geq1$ the hinge is affine on the support and $I(z)=z-\mu$.
Put $z_0=2k+3=2r-1$ and $z_*=1+4r/5$. When $1\leq z\leq z_0$,
the lower-family point $y=(z-3)/2$ gives $M(z)\geq A(z+3)^2/4$ and
\begin{equation}\label{eq:ha-scalar-square}
 \frac{5A}{4}(z+3)^2-(z-\mu)
       =\frac{5A}{4}(z-z_*)^2\geq0.
\end{equation}
For $z\geq z_0$, use $y=k$ and $M(z)\geq b(z-k)$.
The difference $5b(z-k)-(z-\mu)$ is increasing, since $5b-1>0$,
and is nonnegative at $z_0$ by~\eqref{eq:ha-scalar-square}.

For an increasing hinge with $z\geq1$ the claim is trivial. If $z\leq-1$,
then $I(z)=\mu-z$ and $M(z)\geq b(k-z)$. The difference of these bounds
is smallest at $z=-1$, where
$5b(k+1)-(\mu+1)=(25-9r)/10>0$.
For $-1\leq z\leq k$, integration of the lower density gives
\begin{equation}\label{eq:ha-increasing-integral}
 I(z)=\mu-z+b\left((z+1)^2+\frac{(z+1)^3}{6}\right).
\end{equation}
On $[-1,0]$, use $5M(z)\geq5b(k-z)$. The difference from
\eqref{eq:ha-increasing-integral} is concave, with second derivative
$-b(z+3)<0$; its endpoint values are
\[
 \frac{25-9r}{10}>0\quad\text{and}\quad
 \frac{210-91r}{60}>0.
\]
It is therefore positive throughout that interval.
On $[0,k]$, the upper-family point $y=(1+z)/2$ lies in $[k,1]$ and gives
$5M(z)\geq5B(1-z)^2/4$. The difference from
\eqref{eq:ha-increasing-integral} is again concave: its second derivative
is at most $(25-19r)/80<0$. Its endpoint values are
\[
 \frac{555-233r}{480}>0\quad\text{and}\quad
 \frac7{12}B(1-k)^2>0.
\]
Finally, for $k\leq z<1$, only the upper density contributes, and
\[
 I(z)=\frac{C(1-z)^3}{6},\qquad
 M(z)\geq\frac{B(1-z)^2}{4},\qquad
 \frac{I(z)}{M(z)}\leq\frac{2C}{3B}(1-z)\leq\frac83<5,
\]
where $(C/B)(1-k)=4$. This covers all increasing hinges.
\end{proof}

\begin{proof}[Proof of Theorem~\ref{thm:ha-two-hinge}]
Write the selected facet hinge as $f_0=(\alpha\xi+\beta\eta+\gamma)_+$.
Convexity in $\xi$ gives the endpoint bound
\[
 \mathbb E_\xi f_0(\xi,y)
       \leq\frac{f_0(-1,y)+f_0(1,y)}2.
\]
For either edge sign, Lemma~\ref{lem:ha-scalar-amplification} and the four
actual shortest families give
\[
 \mathbb E_\eta f_0(\pm1,\eta)
       \leq5M(f_0(\pm1,\cdot))\leq5R(f_0).
\]
Thus $\mathbb E f_0\leq5R(f_0)$; equation~\eqref{eq:ha-quotient-volume}
proves $R(f_0)\geq2\nu(f_0)$. Symmetry covers every facet.

For two hinges on distinct facet copies, the actual profile is $f+g$.
Nonnegativity yields
\[
 R(f+g)\geq\max\{R(f),R(g)\}
 \geq2\max\{\nu(f),\nu(g)\}\geq\nu(f)+\nu(g).
\]
If both cuts concern the same facet, their actual profile is $h=\max(f,g)$.
Then
\[
 R(h)\geq\max\{R(f),R(g)\}
 \geq2\max\{\nu(f),\nu(g)\}
 \geq\nu(f)+\nu(g)\geq\nu(h).
\]
The simultaneous-splitting derivative theorem gives nonpositive first
variation in either case, with all old rows retained as stipulated.
\end{proof}

The responding shortest word may depend on the hinge. No common occupation
ensemble is asserted. The factor two is what makes these separate responses
sufficient to handle two simultaneous cuts.
